%%%%%%%%%%%%%%%%%%%%%%%%%%%%%%%%%%%%%%%%%%%%%%%%%%%%%%%%%%%%
% 10.6 COMBINATORIAL OPTIMIZATION
% The Composition of Advanced Mathematics
%%%%%%%%%%%%%%%%%%%%%%%%%%%%%%%%%%%%%%%%%%%%%%%%%%%%%%%%%%%%

\section{Combinatorial Optimization}
\label{sec:combinatorial-optimization}

\prerequisite{
Discrete mathematics, graph theory, combinatorics, linear programming,
integer programming, sets, permutations, algorithms, asymptotic
complexity, and basic probability.
}

\learningobjectives{
By the end of this section, the reader should be able to:
\begin{itemize}
    \item define a combinatorial optimization problem;
    \item recognize optimization over subsets, permutations, paths,
          trees, matchings, and discrete structures;
    \item distinguish explicit and implicit feasible sets;
    \item understand why enumeration becomes computationally impossible;
    \item formulate discrete problems using binary variables;
    \item recognize graph-based optimization models;
    \item understand shortest-path optimization;
    \item understand spanning-tree optimization;
    \item understand matching and assignment problems;
    \item formulate the traveling salesperson problem;
    \item understand subtour-elimination constraints;
    \item formulate set-covering and set-packing problems;
    \item understand independent-set, clique, and vertex-cover models;
    \item recognize coloring as a combinatorial optimization problem;
    \item understand knapsack as a combinatorial problem;
    \item recognize matroid optimization;
    \item understand the greedy algorithm for weighted matroid bases;
    \item distinguish exact algorithms from heuristics;
    \item understand dynamic-programming approaches conceptually;
    \item understand branch-and-bound approaches;
    \item understand cutting-plane and branch-and-cut approaches;
    \item use relaxations to obtain bounds;
    \item recognize polyhedral formulations;
    \item understand integrality and total unimodularity conceptually;
    \item distinguish polynomial-time and NP-hard optimization problems;
    \item understand decision versus optimization formulations;
    \item recognize reductions between combinatorial problems;
    \item define approximation algorithms;
    \item understand approximation ratios;
    \item recognize greedy approximation for set cover;
    \item understand metric TSP approximation conceptually;
    \item understand local-search methods;
    \item recognize randomized algorithms conceptually;
    \item understand submodular optimization;
    \item recognize submodular maximization under constraints;
    \item understand decomposition and problem structure;
    \item connect combinatorial optimization with networks, scheduling,
          logistics, computer science, and operations research.
\end{itemize}
}

%%%%%%%%%%%%%%%%%%%%%%%%%%%%%%%%%%%%%%%%%%%%%%%%%%%%%%%%%%%%
\subsection{What Is Combinatorial Optimization?}
\label{subsec:what-is-combinatorial-optimization}
%%%%%%%%%%%%%%%%%%%%%%%%%%%%%%%%%%%%%%%%%%%%%%%%%%%%%%%%%%%%

Combinatorial optimization studies optimization problems whose feasible
solutions are discrete objects.

A general problem can be written

\[
\boxed{
\min_{x\in\mathcal{F}}
c(x),
}
\]

where

\[
\mathcal{F}
\]

is a finite or discrete family of feasible structures.

Examples include:

\[
\boxed{
\text{subsets},
}
\]

\[
\boxed{
\text{permutations},
}
\]

\[
\boxed{
\text{paths},
}
\]

\[
\boxed{
\text{trees},
}
\]

\[
\boxed{
\text{matchings},
}
\]

and

\[
\boxed{
\text{partitions}.
}
\]

%%%%%%%%%%%%%%%%%%%%%%%%%%%%%%%%%%%%%%%%%%%%%%%%%%%%%%%%%%%%
\subsection{Discrete Feasible Structures}
\label{subsec:discrete-feasible-structures}
%%%%%%%%%%%%%%%%%%%%%%%%%%%%%%%%%%%%%%%%%%%%%%%%%%%%%%%%%%%%

Unlike continuous optimization, a combinatorial problem often chooses
among distinct structures.

For example, if \(n\) items may each be selected or not selected, then
there are

\[
\boxed{
2^n
}
\]

possible subsets.

If \(n\) objects may be arranged in any order, there are

\[
\boxed{
n!
}
\]

possible permutations.

These quantities grow extremely rapidly.

%%%%%%%%%%%%%%%%%%%%%%%%%%%%%%%%%%%%%%%%%%%%%%%%%%%%%%%%%%%%
\subsection{Combinatorial Explosion}
\label{subsec:combinatorial-explosion}
%%%%%%%%%%%%%%%%%%%%%%%%%%%%%%%%%%%%%%%%%%%%%%%%%%%%%%%%%%%%

Suppose there are \(50\) binary choices.

Then the number of possible subsets is

\[
\boxed{
2^{50}.
}
\]

This is more than

\[
10^{15}.
\]

Thus complete enumeration can be computationally impossible even for
moderate problem sizes.

Combinatorial optimization seeks mathematical structure that avoids
testing every possible solution.

%%%%%%%%%%%%%%%%%%%%%%%%%%%%%%%%%%%%%%%%%%%%%%%%%%%%%%%%%%%%
\subsection{Explicit and Implicit Feasible Sets}
\label{subsec:explicit-implicit-feasible-sets}
%%%%%%%%%%%%%%%%%%%%%%%%%%%%%%%%%%%%%%%%%%%%%%%%%%%%%%%%%%%%

A feasible set may be listed explicitly:

\[
\mathcal{F}
=
\{
x^{(1)},\ldots,x^{(N)}
\}.
\]

More commonly, the feasible set is described implicitly through rules or
constraints.

For example,

\[
\boxed{
\mathcal{F}
=
\{
S\subseteq V:
|S|=k
\}
}
\]

contains all \(k\)-element subsets of \(V\).

The number of feasible objects may be enormous even when the mathematical
description is short.

%%%%%%%%%%%%%%%%%%%%%%%%%%%%%%%%%%%%%%%%%%%%%%%%%%%%%%%%%%%%
\subsection{Binary Vector Representation}
\label{subsec:binary-vector-combinatorial}
%%%%%%%%%%%%%%%%%%%%%%%%%%%%%%%%%%%%%%%%%%%%%%%%%%%%%%%%%%%%

A subset

\[
S\subseteq\{1,\ldots,n\}
\]

can be represented by its incidence vector

\[
\boxed{
x_i
=
\begin{cases}
1,
&
i\in S,
\\
0,
&
i\notin S.
\end{cases}
}
\]

Thus combinatorial problems can often be written as integer programs.

For example,

\[
|S|\leq k
\]

becomes

\[
\boxed{
\sum_{i=1}^{n}
x_i
\leq
k.
}
\]

%%%%%%%%%%%%%%%%%%%%%%%%%%%%%%%%%%%%%%%%%%%%%%%%%%%%%%%%%%%%
\subsection{Combinatorial Optimization on Graphs}
\label{subsec:combinatorial-graph-optimization}
%%%%%%%%%%%%%%%%%%%%%%%%%%%%%%%%%%%%%%%%%%%%%%%%%%%%%%%%%%%%

Many combinatorial optimization problems are naturally defined on a graph

\[
\boxed{
G=(V,E),
}
\]

where:

\[
V
=
\text{vertex set},
\]

and

\[
E
=
\text{edge set}.
\]

Possible optimization objects include:

\[
\boxed{
\text{paths},
}
\]

\[
\boxed{
\text{cycles},
}
\]

\[
\boxed{
\text{trees},
}
\]

\[
\boxed{
\text{matchings},
}
\]

and

\[
\boxed{
\text{cuts}.
}
\]

%%%%%%%%%%%%%%%%%%%%%%%%%%%%%%%%%%%%%%%%%%%%%%%%%%%%%%%%%%%%
\subsection{Weighted Graphs}
\label{subsec:weighted-graphs-combinatorial}
%%%%%%%%%%%%%%%%%%%%%%%%%%%%%%%%%%%%%%%%%%%%%%%%%%%%%%%%%%%%

Assign a cost

\[
c_e
\]

to each edge

\[
e\in E.
\]

For a selected edge set

\[
F\subseteq E,
\]

its total cost is

\[
\boxed{
c(F)
=
\sum_{e\in F}
c_e.
}
\]

Many graph optimization problems minimize or maximize such sums subject to
structural restrictions on \(F\).

%%%%%%%%%%%%%%%%%%%%%%%%%%%%%%%%%%%%%%%%%%%%%%%%%%%%%%%%%%%%
\subsection{Shortest-Path Problem}
\label{subsec:combinatorial-shortest-path}
%%%%%%%%%%%%%%%%%%%%%%%%%%%%%%%%%%%%%%%%%%%%%%%%%%%%%%%%%%%%

Given a weighted graph and vertices \(s\) and \(t\), the shortest-path
problem asks for an \(s\)-to-\(t\) path of minimum total edge cost.

If

\[
P
\]

is a path, then its cost is

\[
\boxed{
c(P)
=
\sum_{e\in P}
c_e.
}
\]

The optimization problem is

\[
\boxed{
\min_{P\in\mathcal{P}_{st}}
\sum_{e\in P}
c_e.
}
\]

%%%%%%%%%%%%%%%%%%%%%%%%%%%%%%%%%%%%%%%%%%%%%%%%%%%%%%%%%%%%
\subsection{Shortest-Path Algorithms}
\label{subsec:shortest-path-algorithms-combinatorial}
%%%%%%%%%%%%%%%%%%%%%%%%%%%%%%%%%%%%%%%%%%%%%%%%%%%%%%%%%%%%

Important exact algorithms include:

\[
\boxed{
\text{Dijkstra's algorithm}
}
\]

for nonnegative edge weights,

and

\[
\boxed{
\text{Bellman--Ford}
}
\]

when negative edge weights may occur but reachable negative cycles must be
handled appropriately.

Shortest-path problems are examples of combinatorial problems that can be
solved efficiently in polynomial time.

%%%%%%%%%%%%%%%%%%%%%%%%%%%%%%%%%%%%%%%%%%%%%%%%%%%%%%%%%%%%
\subsection{Worked Example: Shortest Path}
\label{subsec:worked-shortest-path-combinatorial}
%%%%%%%%%%%%%%%%%%%%%%%%%%%%%%%%%%%%%%%%%%%%%%%%%%%%%%%%%%%%

\begin{workedexamplebox}{Comparing Candidate Paths}

Suppose three routes from \(A\) to \(D\) have costs

\[
A\to B\to D:
\qquad
3+4=7,
\]

\[
A\to C\to D:
\qquad
2+8=10,
\]

and

\[
A\to B\to C\to D:
\qquad
3+1+2=6.
\]

Therefore,

\begin{answerbox}
\[
\boxed{
A\to B\to C\to D
}
\]

is the shortest among these routes, with cost

\[
\boxed{
6.
}
\]
\end{answerbox}

\end{workedexamplebox}

%%%%%%%%%%%%%%%%%%%%%%%%%%%%%%%%%%%%%%%%%%%%%%%%%%%%%%%%%%%%
\subsection{Minimum Spanning Tree}
\label{subsec:minimum-spanning-tree}
%%%%%%%%%%%%%%%%%%%%%%%%%%%%%%%%%%%%%%%%%%%%%%%%%%%%%%%%%%%%

A spanning tree of a connected graph includes every vertex and contains
no cycles.

A minimum spanning tree, abbreviated MST, solves

\[
\boxed{
\min_{T}
\sum_{e\in T}
c_e
}
\]

over all spanning trees \(T\).

%%%%%%%%%%%%%%%%%%%%%%%%%%%%%%%%%%%%%%%%%%%%%%%%%%%%%%%%%%%%
\subsection{Properties of Trees}
\label{subsec:tree-properties-combinatorial}
%%%%%%%%%%%%%%%%%%%%%%%%%%%%%%%%%%%%%%%%%%%%%%%%%%%%%%%%%%%%

A tree on \(n\) vertices has exactly

\[
\boxed{
n-1
}
\]

edges.

A connected graph with \(n-1\) edges is a tree if and only if it has no
cycles.

These structural properties are fundamental to spanning-tree algorithms.

%%%%%%%%%%%%%%%%%%%%%%%%%%%%%%%%%%%%%%%%%%%%%%%%%%%%%%%%%%%%
\subsection{Kruskal's Algorithm}
\label{subsec:kruskals-algorithm}
%%%%%%%%%%%%%%%%%%%%%%%%%%%%%%%%%%%%%%%%%%%%%%%%%%%%%%%%%%%%

Kruskal's algorithm:

\begin{enumerate}
    \item sorts edges by nondecreasing weight;
    \item scans the edges in that order;
    \item adds an edge if doing so does not create a cycle;
    \item stops after \(n-1\) edges have been selected.
\end{enumerate}

The result is a minimum spanning tree.

%%%%%%%%%%%%%%%%%%%%%%%%%%%%%%%%%%%%%%%%%%%%%%%%%%%%%%%%%%%%
\subsection{Prim's Algorithm}
\label{subsec:prims-algorithm}
%%%%%%%%%%%%%%%%%%%%%%%%%%%%%%%%%%%%%%%%%%%%%%%%%%%%%%%%%%%%

Prim's algorithm grows one connected tree.

Starting from a vertex, repeatedly choose a minimum-cost edge connecting
the current tree to a new vertex.

The algorithm continues until every vertex is included.

%%%%%%%%%%%%%%%%%%%%%%%%%%%%%%%%%%%%%%%%%%%%%%%%%%%%%%%%%%%%
\subsection{Cut Property for MSTs}
\label{subsec:mst-cut-property}
%%%%%%%%%%%%%%%%%%%%%%%%%%%%%%%%%%%%%%%%%%%%%%%%%%%%%%%%%%%%

Suppose a cut partitions the vertices into two sets.

A minimum-weight edge crossing the cut is safe for some minimum spanning
tree under the standard cut-property conditions.

This property provides the theoretical basis for greedy MST algorithms.

%%%%%%%%%%%%%%%%%%%%%%%%%%%%%%%%%%%%%%%%%%%%%%%%%%%%%%%%%%%%
\subsection{Greedy Algorithms}
\label{subsec:greedy-combinatorial}
%%%%%%%%%%%%%%%%%%%%%%%%%%%%%%%%%%%%%%%%%%%%%%%%%%%%%%%%%%%%

A greedy algorithm repeatedly makes a locally attractive choice and never
reverses it.

A generic structure is:

\begin{enumerate}
    \item identify the locally best available choice;
    \item accept it if feasibility is preserved;
    \item repeat until a complete solution is formed.
\end{enumerate}

Greedy methods can be optimal for some combinatorial structures and fail
badly for others.

%%%%%%%%%%%%%%%%%%%%%%%%%%%%%%%%%%%%%%%%%%%%%%%%%%%%%%%%%%%%
\subsection{Why Greedy Methods Need Proof}
\label{subsec:why-greedy-needs-proof}
%%%%%%%%%%%%%%%%%%%%%%%%%%%%%%%%%%%%%%%%%%%%%%%%%%%%%%%%%%%%

A locally best choice need not lead to a globally optimal solution.

Therefore a greedy algorithm requires a structural argument such as:

\[
\boxed{
\text{exchange property},
}
\]

\[
\boxed{
\text{cut property},
}
\]

or

\[
\boxed{
\text{matroid structure}.
}
\]

Without such structure, greedy behavior is only heuristic.

%%%%%%%%%%%%%%%%%%%%%%%%%%%%%%%%%%%%%%%%%%%%%%%%%%%%%%%%%%%%
\subsection{Matching}
\label{subsec:matching-combinatorial}
%%%%%%%%%%%%%%%%%%%%%%%%%%%%%%%%%%%%%%%%%%%%%%%%%%%%%%%%%%%%

A matching is a set of edges such that no two selected edges share a
vertex.

A maximum-cardinality matching solves

\[
\boxed{
\max
|M|
}
\]

over all matchings \(M\).

A weighted matching solves

\[
\boxed{
\max_{M}
\sum_{e\in M}
w_e.
}
\]

%%%%%%%%%%%%%%%%%%%%%%%%%%%%%%%%%%%%%%%%%%%%%%%%%%%%%%%%%%%%
\subsection{Bipartite Matching}
\label{subsec:bipartite-matching}
%%%%%%%%%%%%%%%%%%%%%%%%%%%%%%%%%%%%%%%%%%%%%%%%%%%%%%%%%%%%

A graph is bipartite if

\[
V
=
U\cup W,
\]

where every edge joins a vertex in \(U\) to one in \(W\).

Bipartite matching models:

\[
\boxed{
\text{workers to jobs},
}
\]

\[
\boxed{
\text{students to projects},
}
\]

and

\[
\boxed{
\text{resources to tasks}.
}
\]

%%%%%%%%%%%%%%%%%%%%%%%%%%%%%%%%%%%%%%%%%%%%%%%%%%%%%%%%%%%%
\subsection{Assignment as Matching}
\label{subsec:assignment-as-matching}
%%%%%%%%%%%%%%%%%%%%%%%%%%%%%%%%%%%%%%%%%%%%%%%%%%%%%%%%%%%%

The assignment problem is a weighted perfect matching problem in a
bipartite graph.

Let

\[
x_{ij}\in\{0,1\}
\]

indicate whether agent \(i\) receives task \(j\).

Then

\[
\boxed{
\sum_jx_{ij}=1
}
\]

for each agent,

and

\[
\boxed{
\sum_ix_{ij}=1
}
\]

for each task.

%%%%%%%%%%%%%%%%%%%%%%%%%%%%%%%%%%%%%%%%%%%%%%%%%%%%%%%%%%%%
\subsection{Hungarian Algorithm Preview}
\label{subsec:hungarian-algorithm-preview}
%%%%%%%%%%%%%%%%%%%%%%%%%%%%%%%%%%%%%%%%%%%%%%%%%%%%%%%%%%%%

The assignment problem has specialized polynomial-time algorithms.

One classical approach is the Hungarian algorithm.

Its existence illustrates an important theme:

\[
\boxed{
\text{special combinatorial structure}
\rightarrow
\text{special efficient algorithm}.
}
\]

%%%%%%%%%%%%%%%%%%%%%%%%%%%%%%%%%%%%%%%%%%%%%%%%%%%%%%%%%%%%
\subsection{Hall's Marriage Theorem Connection}
\label{subsec:hall-matching-connection}
%%%%%%%%%%%%%%%%%%%%%%%%%%%%%%%%%%%%%%%%%%%%%%%%%%%%%%%%%%%%

For a bipartite graph with parts \(U\) and \(W\), Hall's theorem gives a
criterion for a matching that covers every vertex of \(U\).

For every subset

\[
S\subseteq U,
\]

the neighborhood must satisfy

\[
\boxed{
|N(S)|
\geq
|S|.
}
\]

This is a structural existence theorem rather than an optimization
algorithm, but it is closely connected with matching theory.

%%%%%%%%%%%%%%%%%%%%%%%%%%%%%%%%%%%%%%%%%%%%%%%%%%%%%%%%%%%%
\subsection{Traveling Salesperson Problem}
\label{subsec:traveling-salesperson-problem}
%%%%%%%%%%%%%%%%%%%%%%%%%%%%%%%%%%%%%%%%%%%%%%%%%%%%%%%%%%%%

The traveling salesperson problem, abbreviated TSP, asks for a
minimum-cost tour that visits every vertex exactly once and returns to the
starting point.

If edge \((i,j)\) has cost \(c_{ij}\), define

\[
x_{ij}
=
\begin{cases}
1,
&
\text{edge }(i,j)\text{ is used},
\\
0,
&
\text{otherwise}.
\end{cases}
\]

The objective is

\[
\boxed{
\min
\sum_{i}
\sum_{j}
c_{ij}x_{ij}.
}
\]

%%%%%%%%%%%%%%%%%%%%%%%%%%%%%%%%%%%%%%%%%%%%%%%%%%%%%%%%%%%%
\subsection{Degree Constraints for TSP}
\label{subsec:tsp-degree-constraints}
%%%%%%%%%%%%%%%%%%%%%%%%%%%%%%%%%%%%%%%%%%%%%%%%%%%%%%%%%%%%

For a directed formulation, each vertex must have one outgoing and one
incoming selected arc:

\[
\boxed{
\sum_{j\neq i}
x_{ij}
=
1,
}
\]

and

\[
\boxed{
\sum_{j\neq i}
x_{ji}
=
1.
}
\]

These conditions alone do not guarantee a single tour.

They can produce several disconnected cycles.

%%%%%%%%%%%%%%%%%%%%%%%%%%%%%%%%%%%%%%%%%%%%%%%%%%%%%%%%%%%%
\subsection{Subtours}
\label{subsec:tsp-subtours}
%%%%%%%%%%%%%%%%%%%%%%%%%%%%%%%%%%%%%%%%%%%%%%%%%%%%%%%%%%%%

A subtour is a cycle containing only a proper subset of the vertices.

For example, the selected arcs might create:

\[
\boxed{
1\to2\to1
}
\]

and separately

\[
\boxed{
3\to4\to5\to3.
}
\]

Every vertex has the correct degree, but there is no single Hamiltonian
tour.

%%%%%%%%%%%%%%%%%%%%%%%%%%%%%%%%%%%%%%%%%%%%%%%%%%%%%%%%%%%%
\subsection{Subtour-Elimination Constraints}
\label{subsec:subtour-elimination}
%%%%%%%%%%%%%%%%%%%%%%%%%%%%%%%%%%%%%%%%%%%%%%%%%%%%%%%%%%%%

For a directed TSP, one family of subtour-elimination constraints is

\[
\boxed{
\sum_{i\in S}
\sum_{j\in S}
x_{ij}
\leq
|S|-1
}
\]

for every nonempty proper subset

\[
S\subsetneq V
\]

with at least two vertices.

These inequalities prevent a closed tour contained entirely inside
\(S\).

%%%%%%%%%%%%%%%%%%%%%%%%%%%%%%%%%%%%%%%%%%%%%%%%%%%%%%%%%%%%
\subsection{Exponential Number of TSP Constraints}
\label{subsec:tsp-exponential-constraints}
%%%%%%%%%%%%%%%%%%%%%%%%%%%%%%%%%%%%%%%%%%%%%%%%%%%%%%%%%%%%

There are exponentially many possible subsets \(S\).

Therefore listing every subtour-elimination inequality in advance can be
impractical.

Instead, modern approaches often:

\[
\boxed{
\text{solve a relaxation},
}
\]

\[
\boxed{
\text{detect violated subtours},
}
\]

and

\[
\boxed{
\text{add cuts as needed}.
}
\]

This is a cutting-plane strategy.

%%%%%%%%%%%%%%%%%%%%%%%%%%%%%%%%%%%%%%%%%%%%%%%%%%%%%%%%%%%%
\subsection{Worked Example: Detecting a Subtour}
\label{subsec:worked-detecting-subtour}
%%%%%%%%%%%%%%%%%%%%%%%%%%%%%%%%%%%%%%%%%%%%%%%%%%%%%%%%%%%%

\begin{workedexamplebox}{A Disconnected TSP Solution}

Suppose a candidate solution contains the cycle

\[
1\to2\to3\to1.
\]

Let

\[
S=\{1,2,3\}.
\]

The three selected arcs lie entirely inside \(S\), so

\[
\sum_{i\in S}
\sum_{j\in S}
x_{ij}
=
3.
\]

But the subtour inequality requires

\[
\sum_{i\in S}
\sum_{j\in S}
x_{ij}
\leq
|S|-1
=
2.
\]

Thus the candidate violates the cut.

\begin{answerbox}
\[
\boxed{
\text{The cycle on }\{1,2,3\}\text{ is eliminated.}
}
\]
\end{answerbox}

\end{workedexamplebox}

%%%%%%%%%%%%%%%%%%%%%%%%%%%%%%%%%%%%%%%%%%%%%%%%%%%%%%%%%%%%
\subsection{Hamiltonian Problems}
\label{subsec:hamiltonian-problems}
%%%%%%%%%%%%%%%%%%%%%%%%%%%%%%%%%%%%%%%%%%%%%%%%%%%%%%%%%%%%

A Hamiltonian cycle visits every vertex exactly once before returning to
the start.

Finding whether one exists is already computationally difficult in
general.

The TSP adds costs and asks for the best such cycle.

Thus optimization can be harder than merely checking feasibility.

%%%%%%%%%%%%%%%%%%%%%%%%%%%%%%%%%%%%%%%%%%%%%%%%%%%%%%%%%%%%
\subsection{Set Cover}
\label{subsec:set-cover}
%%%%%%%%%%%%%%%%%%%%%%%%%%%%%%%%%%%%%%%%%%%%%%%%%%%%%%%%%%%%

Let the universe be

\[
U
=
\{1,\ldots,m\}.
\]

Suppose subsets

\[
S_1,\ldots,S_n
\subseteq U
\]

are available.

The set-cover problem chooses subsets whose union covers \(U\).

Let

\[
x_j\in\{0,1\}
\]

indicate whether \(S_j\) is chosen.

Then

\[
\boxed{
\min
\sum_{j=1}^{n}
c_jx_j
}
\]

subject to

\[
\boxed{
\sum_{j:i\in S_j}
x_j
\geq
1
}
\]

for each element

\[
i\in U.
\]

%%%%%%%%%%%%%%%%%%%%%%%%%%%%%%%%%%%%%%%%%%%%%%%%%%%%%%%%%%%%
\subsection{Set Packing}
\label{subsec:set-packing}
%%%%%%%%%%%%%%%%%%%%%%%%%%%%%%%%%%%%%%%%%%%%%%%%%%%%%%%%%%%%

Set packing chooses a collection of mutually compatible subsets.

A common binary formulation is

\[
\boxed{
\max
\sum_j
w_jx_j
}
\]

subject to

\[
\boxed{
\sum_{j:i\in S_j}
x_j
\leq
1
}
\]

for each underlying element \(i\).

Thus no element may appear in more than one selected set.

%%%%%%%%%%%%%%%%%%%%%%%%%%%%%%%%%%%%%%%%%%%%%%%%%%%%%%%%%%%%
\subsection{Set Partitioning}
\label{subsec:set-partitioning}
%%%%%%%%%%%%%%%%%%%%%%%%%%%%%%%%%%%%%%%%%%%%%%%%%%%%%%%%%%%%

Set partitioning requires each element to be covered exactly once:

\[
\boxed{
\sum_{j:i\in S_j}
x_j
=
1.
}
\]

This structure appears in:

\[
\boxed{
\text{crew scheduling},
}
\]

\[
\boxed{
\text{vehicle routing},
}
\]

and

\[
\boxed{
\text{partitioning problems}.
}
\]

%%%%%%%%%%%%%%%%%%%%%%%%%%%%%%%%%%%%%%%%%%%%%%%%%%%%%%%%%%%%
\subsection{Vertex Cover}
\label{subsec:vertex-cover}
%%%%%%%%%%%%%%%%%%%%%%%%%%%%%%%%%%%%%%%%%%%%%%%%%%%%%%%%%%%%

A vertex cover is a set of vertices touching every edge.

Let

\[
x_v\in\{0,1\}
\]

indicate whether vertex \(v\) is chosen.

A minimum weighted vertex-cover model is

\[
\boxed{
\min
\sum_{v\in V}
w_vx_v
}
\]

subject to

\[
\boxed{
x_u+x_v\geq1
}
\]

for every edge

\[
(u,v)\in E.
\]

%%%%%%%%%%%%%%%%%%%%%%%%%%%%%%%%%%%%%%%%%%%%%%%%%%%%%%%%%%%%
\subsection{Independent Set}
\label{subsec:independent-set-optimization}
%%%%%%%%%%%%%%%%%%%%%%%%%%%%%%%%%%%%%%%%%%%%%%%%%%%%%%%%%%%%

An independent set contains no adjacent vertices.

Let

\[
x_v\in\{0,1\}.
\]

Then

\[
\boxed{
x_u+x_v\leq1
}
\]

for every edge

\[
(u,v)\in E.
\]

The maximum weighted independent-set problem is

\[
\boxed{
\max
\sum_{v\in V}
w_vx_v.
}
\]

%%%%%%%%%%%%%%%%%%%%%%%%%%%%%%%%%%%%%%%%%%%%%%%%%%%%%%%%%%%%
\subsection{Vertex Cover and Independent Set}
\label{subsec:vertex-cover-independent-set}
%%%%%%%%%%%%%%%%%%%%%%%%%%%%%%%%%%%%%%%%%%%%%%%%%%%%%%%%%%%%

For an unweighted graph,

\[
C
\subseteq V
\]

is a vertex cover if and only if

\[
V\setminus C
\]

is an independent set.

Therefore,

\[
\boxed{
\tau(G)+\alpha(G)=|V|,
}
\]

where

\[
\tau(G)
=
\text{minimum vertex-cover size},
\]

and

\[
\alpha(G)
=
\text{maximum independent-set size}.
\]

%%%%%%%%%%%%%%%%%%%%%%%%%%%%%%%%%%%%%%%%%%%%%%%%%%%%%%%%%%%%
\subsection{Clique}
\label{subsec:clique-optimization}
%%%%%%%%%%%%%%%%%%%%%%%%%%%%%%%%%%%%%%%%%%%%%%%%%%%%%%%%%%%%

A clique is a set of pairwise adjacent vertices.

A clique in graph \(G\) is an independent set in the complement graph

\[
\overline{G}.
\]

Therefore,

\[
\boxed{
\omega(G)
=
\alpha(\overline{G}),
}
\]

where \(\omega(G)\) is the maximum clique size.

This gives a direct relationship among several difficult graph
optimization problems.

%%%%%%%%%%%%%%%%%%%%%%%%%%%%%%%%%%%%%%%%%%%%%%%%%%%%%%%%%%%%
\subsection{Graph Coloring}
\label{subsec:graph-coloring-optimization}
%%%%%%%%%%%%%%%%%%%%%%%%%%%%%%%%%%%%%%%%%%%%%%%%%%%%%%%%%%%%

Graph coloring assigns colors to vertices so adjacent vertices receive
different colors.

The optimization problem seeks the minimum number of colors:

\[
\boxed{
\chi(G)
=
\text{chromatic number of }G.
}
\]

Coloring models:

\[
\boxed{
\text{scheduling conflicts},
}
\]

\[
\boxed{
\text{frequency assignment},
}
\]

and

\[
\boxed{
\text{resource incompatibilities}.
}
\]

%%%%%%%%%%%%%%%%%%%%%%%%%%%%%%%%%%%%%%%%%%%%%%%%%%%%%%%%%%%%
\subsection{Binary Coloring Formulation}
\label{subsec:binary-coloring-formulation}
%%%%%%%%%%%%%%%%%%%%%%%%%%%%%%%%%%%%%%%%%%%%%%%%%%%%%%%%%%%%

Let

\[
x_{vc}
=
1
\]

if vertex \(v\) receives color \(c\).

Then every vertex receives exactly one color:

\[
\boxed{
\sum_c
x_{vc}
=
1.
}
\]

For every edge \((u,v)\),

\[
\boxed{
x_{uc}+x_{vc}
\leq1
}
\]

for each color \(c\).

Additional variables can indicate whether a color is used.

%%%%%%%%%%%%%%%%%%%%%%%%%%%%%%%%%%%%%%%%%%%%%%%%%%%%%%%%%%%%
\subsection{Maximum Cut}
\label{subsec:maximum-cut}
%%%%%%%%%%%%%%%%%%%%%%%%%%%%%%%%%%%%%%%%%%%%%%%%%%%%%%%%%%%%

A cut partitions

\[
V
=
S\cup(V\setminus S).
\]

An edge crosses the cut when its endpoints lie on opposite sides.

The maximum-cut problem seeks a partition maximizing the total weight of
crossing edges:

\[
\boxed{
\max
\sum_{e\in\delta(S)}
w_e.
}
\]

This is a fundamental NP-hard combinatorial optimization problem.

%%%%%%%%%%%%%%%%%%%%%%%%%%%%%%%%%%%%%%%%%%%%%%%%%%%%%%%%%%%%
\subsection{Knapsack Revisited}
\label{subsec:knapsack-combinatorial}
%%%%%%%%%%%%%%%%%%%%%%%%%%%%%%%%%%%%%%%%%%%%%%%%%%%%%%%%%%%%

The \(0\)-\(1\) knapsack problem is

\[
\boxed{
\begin{aligned}
\max
\quad&
\sum_j
v_jx_j
\\
\text{subject to}
\quad&
\sum_j
w_jx_j
\leq
W,
\\
&
x_j\in\{0,1\}.
\end{aligned}
}
\]

It illustrates several major combinatorial ideas:

\[
\boxed{
\text{binary modeling},
}
\]

\[
\boxed{
\text{dynamic programming},
}
\]

\[
\boxed{
\text{branch-and-bound},
}
\]

and

\[
\boxed{
\text{approximation algorithms}.
}
\]

%%%%%%%%%%%%%%%%%%%%%%%%%%%%%%%%%%%%%%%%%%%%%%%%%%%%%%%%%%%%
\subsection{Matroids}
\label{subsec:matroids-combinatorial-optimization}
%%%%%%%%%%%%%%%%%%%%%%%%%%%%%%%%%%%%%%%%%%%%%%%%%%%%%%%%%%%%

A matroid is a pair

\[
\boxed{
M=(E,\mathcal{I}),
}
\]

where \(E\) is a finite ground set and \(\mathcal{I}\) is a family of
independent subsets satisfying matroid axioms.

Matroids abstract independence from:

\[
\boxed{
\text{linear algebra},
}
\]

and

\[
\boxed{
\text{graph forests}.
}
\]

They provide a precise setting in which greedy optimization is optimal.

%%%%%%%%%%%%%%%%%%%%%%%%%%%%%%%%%%%%%%%%%%%%%%%%%%%%%%%%%%%%
\subsection{Matroid Independence Axioms}
\label{subsec:matroid-independence-axioms}
%%%%%%%%%%%%%%%%%%%%%%%%%%%%%%%%%%%%%%%%%%%%%%%%%%%%%%%%%%%%

A standard independence formulation requires:

\[
\boxed{
\varnothing\in\mathcal{I},
}
\]

and hereditary closure:

\[
\boxed{
A\in\mathcal{I},
\quad
B\subseteq A
\Rightarrow
B\in\mathcal{I}.
}
\]

The exchange property states that if

\[
A,B\in\mathcal{I}
\]

and

\[
|A|<|B|,
\]

then there exists

\[
e\in B\setminus A
\]

such that

\[
\boxed{
A\cup\{e\}\in\mathcal{I}.
}
\]

%%%%%%%%%%%%%%%%%%%%%%%%%%%%%%%%%%%%%%%%%%%%%%%%%%%%%%%%%%%%
\subsection{Matroid Bases}
\label{subsec:matroid-bases-optimization}
%%%%%%%%%%%%%%%%%%%%%%%%%%%%%%%%%%%%%%%%%%%%%%%%%%%%%%%%%%%%

A basis is a maximal independent set.

All bases of a matroid have the same cardinality.

Given weights

\[
w_e,
\]

the maximum-weight basis problem is

\[
\boxed{
\max_{B\text{ basis}}
\sum_{e\in B}
w_e.
}
\]

%%%%%%%%%%%%%%%%%%%%%%%%%%%%%%%%%%%%%%%%%%%%%%%%%%%%%%%%%%%%
\subsection{Greedy Algorithm for Matroid Bases}
\label{subsec:matroid-greedy-algorithm}
%%%%%%%%%%%%%%%%%%%%%%%%%%%%%%%%%%%%%%%%%%%%%%%%%%%%%%%%%%%%

To find a maximum-weight basis:

\begin{enumerate}
    \item order the elements so that
    \[
    w_{e_1}\geq w_{e_2}\geq\cdots\geq w_{e_m};
    \]
    \item begin with \(S=\varnothing\);
    \item scan elements in that order;
    \item add \(e_i\) whenever
    \[
    S\cup\{e_i\}\in\mathcal{I};
    \]
    \item stop when a basis is obtained.
\end{enumerate}

The matroid exchange property guarantees optimality.

%%%%%%%%%%%%%%%%%%%%%%%%%%%%%%%%%%%%%%%%%%%%%%%%%%%%%%%%%%%%
\subsection{Graphic Matroid}
\label{subsec:graphic-matroid-optimization}
%%%%%%%%%%%%%%%%%%%%%%%%%%%%%%%%%%%%%%%%%%%%%%%%%%%%%%%%%%%%

For a graph \(G=(V,E)\), define an edge set to be independent if it
contains no cycle.

Then

\[
\boxed{
(E,\mathcal{I})
}
\]

forms the graphic matroid.

Its bases are spanning forests, and for a connected graph, spanning
trees.

Thus Kruskal's algorithm is an instance of matroid greedy optimization.

%%%%%%%%%%%%%%%%%%%%%%%%%%%%%%%%%%%%%%%%%%%%%%%%%%%%%%%%%%%%
\subsection{Greedy Failure Outside Matroids}
\label{subsec:greedy-failure-outside-matroids}
%%%%%%%%%%%%%%%%%%%%%%%%%%%%%%%%%%%%%%%%%%%%%%%%%%%%%%%%%%%%

Consider the \(0\)-\(1\) knapsack problem.

Selecting items greedily by largest value does not always work.

Selecting by largest value-to-weight ratio also does not always produce
the optimal \(0\)-\(1\) solution.

Therefore:

\[
\boxed{
\text{greedy success depends on structure}.
}
\]

%%%%%%%%%%%%%%%%%%%%%%%%%%%%%%%%%%%%%%%%%%%%%%%%%%%%%%%%%%%%
\subsection{Exact Algorithms}
\label{subsec:exact-combinatorial-algorithms}
%%%%%%%%%%%%%%%%%%%%%%%%%%%%%%%%%%%%%%%%%%%%%%%%%%%%%%%%%%%%

An exact algorithm guarantees an optimal solution.

Examples include:

\[
\boxed{
\text{specialized polynomial algorithms},
}
\]

\[
\boxed{
\text{dynamic programming},
}
\]

\[
\boxed{
\text{branch-and-bound},
}
\]

and

\[
\boxed{
\text{branch-and-cut}.
}
\]

The computational cost may vary from polynomial to exponential.

%%%%%%%%%%%%%%%%%%%%%%%%%%%%%%%%%%%%%%%%%%%%%%%%%%%%%%%%%%%%
\subsection{Enumeration}
\label{subsec:enumeration-combinatorial}
%%%%%%%%%%%%%%%%%%%%%%%%%%%%%%%%%%%%%%%%%%%%%%%%%%%%%%%%%%%%

The simplest exact method is complete enumeration:

\[
\boxed{
\text{evaluate every feasible solution}.
}
\]

If there are \(2^n\) subsets, this requires exponential work in the worst
case.

Enumeration is therefore useful mainly for:

\[
\boxed{
\text{small instances},
}
\]

or

\[
\boxed{
\text{benchmarking other methods}.
}
\]

%%%%%%%%%%%%%%%%%%%%%%%%%%%%%%%%%%%%%%%%%%%%%%%%%%%%%%%%%%%%
\subsection{Dynamic Programming Preview}
\label{subsec:dynamic-programming-combinatorial-preview}
%%%%%%%%%%%%%%%%%%%%%%%%%%%%%%%%%%%%%%%%%%%%%%%%%%%%%%%%%%%%

Dynamic programming breaks a problem into overlapping subproblems.

It is effective when the problem has:

\[
\boxed{
\text{optimal substructure},
}
\]

and

\[
\boxed{
\text{repeated subproblems}.
}
\]

Instead of solving the same subproblem repeatedly, its solution is stored
and reused.

Dynamic Programming is developed in detail in
Section~\ref{sec:dynamic-programming}.

%%%%%%%%%%%%%%%%%%%%%%%%%%%%%%%%%%%%%%%%%%%%%%%%%%%%%%%%%%%%
\subsection{Branch-and-Bound in Combinatorial Optimization}
\label{subsec:branch-and-bound-combinatorial}
%%%%%%%%%%%%%%%%%%%%%%%%%%%%%%%%%%%%%%%%%%%%%%%%%%%%%%%%%%%%

Branch-and-bound partitions the discrete feasible set into subproblems.

For each subproblem:

\begin{enumerate}
    \item compute a relaxation bound;
    \item compare the bound with the incumbent;
    \item prune if improvement is impossible;
    \item otherwise branch further.
\end{enumerate}

This framework applies far beyond generic integer programming.

%%%%%%%%%%%%%%%%%%%%%%%%%%%%%%%%%%%%%%%%%%%%%%%%%%%%%%%%%%%%
\subsection{Relaxations}
\label{subsec:combinatorial-relaxations}
%%%%%%%%%%%%%%%%%%%%%%%%%%%%%%%%%%%%%%%%%%%%%%%%%%%%%%%%%%%%

A relaxation replaces a difficult feasible set \(\mathcal{F}\) by an
easier superset \(\mathcal{R}\):

\[
\boxed{
\mathcal{F}
\subseteq
\mathcal{R}.
}
\]

For minimization,

\[
\boxed{
\min_{x\in\mathcal{R}}
f(x)
\leq
\min_{x\in\mathcal{F}}
f(x).
}
\]

Thus the relaxation gives a lower bound.

For maximization, it gives an upper bound.

%%%%%%%%%%%%%%%%%%%%%%%%%%%%%%%%%%%%%%%%%%%%%%%%%%%%%%%%%%%%
\subsection{Common Relaxations}
\label{subsec:common-combinatorial-relaxations}
%%%%%%%%%%%%%%%%%%%%%%%%%%%%%%%%%%%%%%%%%%%%%%%%%%%%%%%%%%%%

Important relaxation types include:

\[
\boxed{
\text{linear-programming relaxations},
}
\]

\[
\boxed{
\text{Lagrangian relaxations},
}
\]

\[
\boxed{
\text{semidefinite relaxations},
}
\]

and

\[
\boxed{
\text{problem-specific combinatorial bounds}.
}
\]

A good relaxation should be both strong and computationally manageable.

%%%%%%%%%%%%%%%%%%%%%%%%%%%%%%%%%%%%%%%%%%%%%%%%%%%%%%%%%%%%
\subsection{Polyhedral Approach}
\label{subsec:polyhedral-combinatorial-optimization}
%%%%%%%%%%%%%%%%%%%%%%%%%%%%%%%%%%%%%%%%%%%%%%%%%%%%%%%%%%%%

Represent each feasible combinatorial object by a binary incidence vector.

If

\[
X
\subseteq
\{0,1\}^n
\]

is the set of feasible incidence vectors, study the polytope

\[
\boxed{
P
=
\operatorname{conv}(X).
}
\]

A perfect linear description of \(P\) would allow linear programming to
solve the combinatorial problem.

%%%%%%%%%%%%%%%%%%%%%%%%%%%%%%%%%%%%%%%%%%%%%%%%%%%%%%%%%%%%
\subsection{Facet-Defining Inequalities}
\label{subsec:facet-defining-inequalities}
%%%%%%%%%%%%%%%%%%%%%%%%%%%%%%%%%%%%%%%%%%%%%%%%%%%%%%%%%%%%

Strong linear inequalities defining high-dimensional faces of a
combinatorial polytope are called facet-defining inequalities when they
define facets.

Finding such inequalities is central to polyhedral combinatorics.

Strong facets can substantially improve branch-and-cut algorithms.

%%%%%%%%%%%%%%%%%%%%%%%%%%%%%%%%%%%%%%%%%%%%%%%%%%%%%%%%%%%%
\subsection{Separation Problem}
\label{subsec:combinatorial-separation-problem}
%%%%%%%%%%%%%%%%%%%%%%%%%%%%%%%%%%%%%%%%%%%%%%%%%%%%%%%%%%%%

Given a point \(\bar{x}\), the separation problem asks:

\[
\boxed{
\text{Is }\bar{x}\in P?
}
\]

If not, find a valid inequality

\[
\boxed{
a^Tx\leq b
}
\]

satisfied by every point in \(P\) but violated by \(\bar{x}\).

Separation is the algorithmic counterpart of cutting-plane generation.

%%%%%%%%%%%%%%%%%%%%%%%%%%%%%%%%%%%%%%%%%%%%%%%%%%%%%%%%%%%%
\subsection{Optimization and Separation}
\label{subsec:optimization-separation}
%%%%%%%%%%%%%%%%%%%%%%%%%%%%%%%%%%%%%%%%%%%%%%%%%%%%%%%%%%%%

A deep principle in combinatorial optimization is that, under suitable
conditions, polynomial-time linear optimization over a rational polytope
and polynomial-time separation over that polytope are closely related.

This connects geometric descriptions with algorithms.

The full equivalence is developed in advanced polyhedral optimization.

%%%%%%%%%%%%%%%%%%%%%%%%%%%%%%%%%%%%%%%%%%%%%%%%%%%%%%%%%%%%
\subsection{Integrality}
\label{subsec:combinatorial-integrality}
%%%%%%%%%%%%%%%%%%%%%%%%%%%%%%%%%%%%%%%%%%%%%%%%%%%%%%%%%%%%

A polyhedron is integral if every extreme point is integral.

If the feasible relaxation polyhedron is integral, then

\[
\boxed{
\text{linear programming solves the discrete problem directly}.
}
\]

Examples occur in:

\[
\boxed{
\text{shortest paths},
}
\]

\[
\boxed{
\text{network flows},
}
\]

and

\[
\boxed{
\text{bipartite matching}.
}
\]

%%%%%%%%%%%%%%%%%%%%%%%%%%%%%%%%%%%%%%%%%%%%%%%%%%%%%%%%%%%%
\subsection{Total Unimodularity Connection}
\label{subsec:combinatorial-total-unimodularity}
%%%%%%%%%%%%%%%%%%%%%%%%%%%%%%%%%%%%%%%%%%%%%%%%%%%%%%%%%%%%

Total unimodularity provides one important sufficient condition for
integrality.

If

\[
A
\]

is totally unimodular and

\[
b
\]

is integral, then under appropriate standard formulations the polyhedron

\[
\boxed{
\{
x:
Ax\leq b
\}
}
\]

has integral vertices.

This explains why certain discrete problems have unusually tractable LP
formulations.

%%%%%%%%%%%%%%%%%%%%%%%%%%%%%%%%%%%%%%%%%%%%%%%%%%%%%%%%%%%%
\subsection{Computational Complexity}
\label{subsec:combinatorial-complexity}
%%%%%%%%%%%%%%%%%%%%%%%%%%%%%%%%%%%%%%%%%%%%%%%%%%%%%%%%%%%%

Combinatorial optimization problems differ dramatically in computational
difficulty.

Some admit polynomial-time algorithms.

Others are NP-hard.

Examples of polynomially solvable problems include:

\[
\boxed{
\text{shortest path},
}
\]

\[
\boxed{
\text{minimum spanning tree},
}
\]

and

\[
\boxed{
\text{bipartite matching}.
}
\]

Examples of NP-hard problems include:

\[
\boxed{
\text{general TSP},
}
\]

\[
\boxed{
\text{maximum clique},
}
\]

\[
\boxed{
\text{graph coloring},
}
\]

and

\[
\boxed{
\text{set cover}.
}
\]

%%%%%%%%%%%%%%%%%%%%%%%%%%%%%%%%%%%%%%%%%%%%%%%%%%%%%%%%%%%%
\subsection{Decision Problems}
\label{subsec:decision-problems-combinatorial}
%%%%%%%%%%%%%%%%%%%%%%%%%%%%%%%%%%%%%%%%%%%%%%%%%%%%%%%%%%%%

An optimization problem can often be associated with a yes-or-no decision
problem.

For example, instead of asking for the minimum TSP tour length, ask:

\[
\boxed{
\text{Is there a tour of cost at most }B?
}
\]

Decision versions are central to computational-complexity theory.

%%%%%%%%%%%%%%%%%%%%%%%%%%%%%%%%%%%%%%%%%%%%%%%%%%%%%%%%%%%%
\subsection{Class P}
\label{subsec:class-p-combinatorial}
%%%%%%%%%%%%%%%%%%%%%%%%%%%%%%%%%%%%%%%%%%%%%%%%%%%%%%%%%%%%

The class \(P\) contains decision problems solvable in polynomial time by
a deterministic algorithm.

Informally, these are problems for which an algorithm exists whose running
time is bounded by

\[
\boxed{
O(n^k)
}
\]

for some fixed constant \(k\), measured in the appropriate input size.

%%%%%%%%%%%%%%%%%%%%%%%%%%%%%%%%%%%%%%%%%%%%%%%%%%%%%%%%%%%%
\subsection{Class NP}
\label{subsec:class-np-combinatorial}
%%%%%%%%%%%%%%%%%%%%%%%%%%%%%%%%%%%%%%%%%%%%%%%%%%%%%%%%%%%%

The class \(NP\) contains decision problems for which a proposed
yes-certificate can be verified in polynomial time.

For example, given a proposed TSP tour, one can check:

\[
\boxed{
\text{whether every vertex appears exactly once},
}
\]

and

\[
\boxed{
\text{whether its cost is at most }B
}
\]

in polynomial time.

%%%%%%%%%%%%%%%%%%%%%%%%%%%%%%%%%%%%%%%%%%%%%%%%%%%%%%%%%%%%
\subsection{NP-Complete Problems}
\label{subsec:np-complete-combinatorial}
%%%%%%%%%%%%%%%%%%%%%%%%%%%%%%%%%%%%%%%%%%%%%%%%%%%%%%%%%%%%

A decision problem is NP-complete if:

\begin{enumerate}
    \item it belongs to \(NP\); and
    \item every problem in \(NP\) can be reduced to it in polynomial time.
\end{enumerate}

NP-complete problems are among the hardest problems in \(NP\) under
polynomial-time reductions.

%%%%%%%%%%%%%%%%%%%%%%%%%%%%%%%%%%%%%%%%%%%%%%%%%%%%%%%%%%%%
\subsection{NP-Hard Optimization Problems}
\label{subsec:np-hard-optimization}
%%%%%%%%%%%%%%%%%%%%%%%%%%%%%%%%%%%%%%%%%%%%%%%%%%%%%%%%%%%%

An optimization problem is NP-hard if solving it efficiently would allow
one to solve an NP-hard or NP-complete decision problem efficiently.

NP-hardness does not imply that every instance is difficult.

It describes worst-case computational complexity.

%%%%%%%%%%%%%%%%%%%%%%%%%%%%%%%%%%%%%%%%%%%%%%%%%%%%%%%%%%%%
\subsection{Polynomial-Time Reductions}
\label{subsec:polynomial-time-reductions}
%%%%%%%%%%%%%%%%%%%%%%%%%%%%%%%%%%%%%%%%%%%%%%%%%%%%%%%%%%%%

A polynomial-time reduction transforms instances of problem \(A\) into
instances of problem \(B\) so that solving \(B\) solves \(A\).

Symbolically,

\[
\boxed{
A
\leq_P
B.
}
\]

If \(A\) is known to be hard and

\[
A\leq_P B,
\]

then \(B\) is at least as hard as \(A\) in the reduction sense.

%%%%%%%%%%%%%%%%%%%%%%%%%%%%%%%%%%%%%%%%%%%%%%%%%%%%%%%%%%%%
\subsection{Why Complexity Classification Matters}
\label{subsec:why-complexity-matters}
%%%%%%%%%%%%%%%%%%%%%%%%%%%%%%%%%%%%%%%%%%%%%%%%%%%%%%%%%%%%

Complexity guides algorithm selection.

If a problem is polynomially solvable, seek an exact specialized
algorithm.

If it is NP-hard, one may need:

\[
\boxed{
\text{branch-and-bound},
}
\]

\[
\boxed{
\text{approximations},
}
\]

\[
\boxed{
\text{heuristics},
}
\]

or

\[
\boxed{
\text{problem-specific structure}.
}
\]

%%%%%%%%%%%%%%%%%%%%%%%%%%%%%%%%%%%%%%%%%%%%%%%%%%%%%%%%%%%%
\subsection{Approximation Algorithms}
\label{subsec:approximation-algorithms}
%%%%%%%%%%%%%%%%%%%%%%%%%%%%%%%%%%%%%%%%%%%%%%%%%%%%%%%%%%%%

An approximation algorithm runs efficiently and guarantees a solution
within a known factor of optimal.

For a minimization problem, an algorithm is a \(\rho\)-approximation if

\[
\boxed{
C_{\mathrm{alg}}
\leq
\rho
C^*
}
\]

for every instance, where

\[
\rho\geq1.
\]

%%%%%%%%%%%%%%%%%%%%%%%%%%%%%%%%%%%%%%%%%%%%%%%%%%%%%%%%%%%%
\subsection{Approximation Ratio for Maximization}
\label{subsec:maximization-approximation-ratio}
%%%%%%%%%%%%%%%%%%%%%%%%%%%%%%%%%%%%%%%%%%%%%%%%%%%%%%%%%%%%

For maximization, a common guarantee is

\[
\boxed{
C_{\mathrm{alg}}
\geq
\alpha
C^*,
}
\]

where

\[
0<\alpha\leq1.
\]

Equivalent ratio conventions are also used.

The convention should always be stated.

%%%%%%%%%%%%%%%%%%%%%%%%%%%%%%%%%%%%%%%%%%%%%%%%%%%%%%%%%%%%
\subsection{Approximation Versus Heuristic}
\label{subsec:approximation-versus-heuristic}
%%%%%%%%%%%%%%%%%%%%%%%%%%%%%%%%%%%%%%%%%%%%%%%%%%%%%%%%%%%%

A heuristic aims to find good solutions but may have no proven
worst-case guarantee.

An approximation algorithm provides a theorem such as

\[
\boxed{
C_{\mathrm{alg}}
\leq
2C^*.
}
\]

Thus:

\[
\boxed{
\text{heuristic}
\neq
\text{approximation algorithm}.
}
\]

%%%%%%%%%%%%%%%%%%%%%%%%%%%%%%%%%%%%%%%%%%%%%%%%%%%%%%%%%%%%
\subsection{Greedy Set Cover}
\label{subsec:greedy-set-cover}
%%%%%%%%%%%%%%%%%%%%%%%%%%%%%%%%%%%%%%%%%%%%%%%%%%%%%%%%%%%%

For unweighted set cover, a natural greedy algorithm repeatedly chooses
the set covering the largest number of currently uncovered elements.

For weighted set cover, choose a set with favorable cost per newly covered
element:

\[
\boxed{
\frac{
c_j
}{
\text{number of newly covered elements}
}.
}
\]

%%%%%%%%%%%%%%%%%%%%%%%%%%%%%%%%%%%%%%%%%%%%%%%%%%%%%%%%%%%%
\subsection{Set-Cover Approximation Guarantee}
\label{subsec:set-cover-approximation}
%%%%%%%%%%%%%%%%%%%%%%%%%%%%%%%%%%%%%%%%%%%%%%%%%%%%%%%%%%%%

The greedy set-cover algorithm has a logarithmic approximation guarantee.

For a universe of size \(m\), the approximation factor can be bounded by a
harmonic-number quantity

\[
\boxed{
H_m
=
1+\frac12+\frac13+\cdots+\frac1m,
}
\]

with

\[
\boxed{
H_m
\leq
1+\ln m.
}
\]

%%%%%%%%%%%%%%%%%%%%%%%%%%%%%%%%%%%%%%%%%%%%%%%%%%%%%%%%%%%%
\subsection{Metric TSP}
\label{subsec:metric-tsp}
%%%%%%%%%%%%%%%%%%%%%%%%%%%%%%%%%%%%%%%%%%%%%%%%%%%%%%%%%%%%

A TSP instance is metric if edge costs satisfy:

\[
\boxed{
c_{ij}\geq0,
}
\]

\[
\boxed{
c_{ij}=c_{ji},
}
\]

and the triangle inequality

\[
\boxed{
c_{ij}
\leq
c_{ik}+c_{kj}.
}
\]

Metric structure permits approximation guarantees that are unavailable
for arbitrary TSP costs.

%%%%%%%%%%%%%%%%%%%%%%%%%%%%%%%%%%%%%%%%%%%%%%%%%%%%%%%%%%%%
\subsection{Double-Tree Approximation for Metric TSP}
\label{subsec:double-tree-tsp}
%%%%%%%%%%%%%%%%%%%%%%%%%%%%%%%%%%%%%%%%%%%%%%%%%%%%%%%%%%%%

A basic metric-TSP approximation proceeds as follows:

\begin{enumerate}
    \item compute a minimum spanning tree;
    \item duplicate every tree edge;
    \item form an Eulerian tour;
    \item shortcut repeated vertices.
\end{enumerate}

By the triangle inequality, shortcutting does not increase cost.

The resulting tour has cost at most

\[
\boxed{
2C^*.
}
\]

%%%%%%%%%%%%%%%%%%%%%%%%%%%%%%%%%%%%%%%%%%%%%%%%%%%%%%%%%%%%
\subsection{Why the Double-Tree Bound Works}
\label{subsec:double-tree-bound}
%%%%%%%%%%%%%%%%%%%%%%%%%%%%%%%%%%%%%%%%%%%%%%%%%%%%%%%%%%%%

Removing one edge from an optimal TSP tour produces a spanning tree.

Therefore,

\[
\boxed{
C_{\mathrm{MST}}
\leq
C^*_{\mathrm{TSP}}.
}
\]

Duplicating the MST gives cost

\[
2C_{\mathrm{MST}}
\leq
2C^*_{\mathrm{TSP}}.
\]

Metric shortcutting preserves the upper bound.

%%%%%%%%%%%%%%%%%%%%%%%%%%%%%%%%%%%%%%%%%%%%%%%%%%%%%%%%%%%%
\subsection{Christofides Algorithm Preview}
\label{subsec:christofides-preview}
%%%%%%%%%%%%%%%%%%%%%%%%%%%%%%%%%%%%%%%%%%%%%%%%%%%%%%%%%%%%

For symmetric metric TSP, Christofides' algorithm combines:

\[
\boxed{
\text{minimum spanning tree},
}
\]

\[
\boxed{
\text{minimum-weight matching on odd-degree vertices},
}
\]

and

\[
\boxed{
\text{Eulerian shortcutting}.
}
\]

Its classical approximation factor is

\[
\boxed{
\frac32.
}
\]

%%%%%%%%%%%%%%%%%%%%%%%%%%%%%%%%%%%%%%%%%%%%%%%%%%%%%%%%%%%%
\subsection{Vertex-Cover Approximation Preview}
\label{subsec:vertex-cover-approximation}
%%%%%%%%%%%%%%%%%%%%%%%%%%%%%%%%%%%%%%%%%%%%%%%%%%%%%%%%%%%%

A simple algorithm for unweighted vertex cover chooses both endpoints of
edges in a maximal matching.

If \(M\) is the matching, the resulting cover has size

\[
2|M|.
\]

Every vertex cover must contain at least one endpoint of every matching
edge.

Thus

\[
\boxed{
|C|
\leq
2|C^*|.
}
\]

This gives a \(2\)-approximation.

%%%%%%%%%%%%%%%%%%%%%%%%%%%%%%%%%%%%%%%%%%%%%%%%%%%%%%%%%%%%
\subsection{Local Search}
\label{subsec:local-search-combinatorial}
%%%%%%%%%%%%%%%%%%%%%%%%%%%%%%%%%%%%%%%%%%%%%%%%%%%%%%%%%%%%

Local search begins from a feasible discrete solution and repeatedly moves
to a neighboring solution with improved objective value.

A generic scheme is:

\begin{enumerate}
    \item construct an initial feasible solution;
    \item define a neighborhood;
    \item search for an improving neighbor;
    \item move to that neighbor;
    \item stop when no improving move is found.
\end{enumerate}

%%%%%%%%%%%%%%%%%%%%%%%%%%%%%%%%%%%%%%%%%%%%%%%%%%%%%%%%%%%%
\subsection{Neighborhood Structures}
\label{subsec:combinatorial-neighborhoods}
%%%%%%%%%%%%%%%%%%%%%%%%%%%%%%%%%%%%%%%%%%%%%%%%%%%%%%%%%%%%

Examples include:

\[
\boxed{
\text{swap one selected item},
}
\]

\[
\boxed{
\text{reverse part of a route},
}
\]

\[
\boxed{
\text{move one job between machines},
}
\]

or

\[
\boxed{
\text{exchange two assignments}.
}
\]

Neighborhood design determines what local optima mean.

%%%%%%%%%%%%%%%%%%%%%%%%%%%%%%%%%%%%%%%%%%%%%%%%%%%%%%%%%%%%
\subsection{2-Opt for TSP}
\label{subsec:two-opt-tsp}
%%%%%%%%%%%%%%%%%%%%%%%%%%%%%%%%%%%%%%%%%%%%%%%%%%%%%%%%%%%%

A \(2\)-opt move removes two tour edges and reconnects the resulting paths
in a different way.

If the new tour is shorter, accept the move.

Repeated \(2\)-opt improvement can produce good tours quickly.

However, a \(2\)-opt local optimum need not be globally optimal.

%%%%%%%%%%%%%%%%%%%%%%%%%%%%%%%%%%%%%%%%%%%%%%%%%%%%%%%%%%%%
\subsection{Metaheuristics Preview}
\label{subsec:metaheuristics-preview}
%%%%%%%%%%%%%%%%%%%%%%%%%%%%%%%%%%%%%%%%%%%%%%%%%%%%%%%%%%%%

More advanced heuristic frameworks include:

\[
\boxed{
\text{simulated annealing},
}
\]

\[
\boxed{
\text{tabu search},
}
\]

\[
\boxed{
\text{genetic algorithms},
}
\]

and

\[
\boxed{
\text{variable-neighborhood search}.
}
\]

These methods can be effective in practice but generally do not provide
global optimality certificates.

%%%%%%%%%%%%%%%%%%%%%%%%%%%%%%%%%%%%%%%%%%%%%%%%%%%%%%%%%%%%
\subsection{Randomized Algorithms}
\label{subsec:randomized-combinatorial-algorithms}
%%%%%%%%%%%%%%%%%%%%%%%%%%%%%%%%%%%%%%%%%%%%%%%%%%%%%%%%%%%%

A randomized algorithm makes random choices during computation.

Its performance may be analyzed through:

\[
\boxed{
\text{expected running time},
}
\]

\[
\boxed{
\text{probability of success},
}
\]

or

\[
\boxed{
\text{expected approximation quality}.
}
\]

Randomization can simplify algorithms or improve practical performance.

%%%%%%%%%%%%%%%%%%%%%%%%%%%%%%%%%%%%%%%%%%%%%%%%%%%%%%%%%%%%
\subsection{Randomized Rounding Preview}
\label{subsec:randomized-rounding-preview}
%%%%%%%%%%%%%%%%%%%%%%%%%%%%%%%%%%%%%%%%%%%%%%%%%%%%%%%%%%%%

Suppose an LP relaxation gives fractional values

\[
0\leq x_i^*\leq1.
\]

Randomized rounding may set

\[
X_i
=
1
\]

with probability

\[
\boxed{
x_i^*.
}
\]

Then

\[
\boxed{
\mathbb{E}[X_i]
=
x_i^*.
}
\]

The challenge is controlling feasibility and concentration around expected
values.

%%%%%%%%%%%%%%%%%%%%%%%%%%%%%%%%%%%%%%%%%%%%%%%%%%%%%%%%%%%%
\subsection{Probabilistic Method Connection}
\label{subsec:probabilistic-method-combinatorial}
%%%%%%%%%%%%%%%%%%%%%%%%%%%%%%%%%%%%%%%%%%%%%%%%%%%%%%%%%%%%

Sometimes random selection proves that a good combinatorial object must
exist.

If a randomly generated solution has positive probability of satisfying
a desired property, then at least one such deterministic object exists.

This connects combinatorial optimization with probabilistic combinatorics.

%%%%%%%%%%%%%%%%%%%%%%%%%%%%%%%%%%%%%%%%%%%%%%%%%%%%%%%%%%%%
\subsection{Submodular Functions}
\label{subsec:submodular-functions-optimization}
%%%%%%%%%%%%%%%%%%%%%%%%%%%%%%%%%%%%%%%%%%%%%%%%%%%%%%%%%%%%

A set function

\[
f:2^V\to\mathbb{R}
\]

is submodular if it has diminishing returns.

One equivalent condition is

\[
\boxed{
f(A\cup\{e\})-f(A)
\geq
f(B\cup\{e\})-f(B)
}
\]

whenever

\[
A\subseteq B
\]

and

\[
e\notin B.
\]

The marginal value of adding \(e\) decreases as the selected set grows.

%%%%%%%%%%%%%%%%%%%%%%%%%%%%%%%%%%%%%%%%%%%%%%%%%%%%%%%%%%%%
\subsection{Examples of Submodularity}
\label{subsec:examples-submodularity}
%%%%%%%%%%%%%%%%%%%%%%%%%%%%%%%%%%%%%%%%%%%%%%%%%%%%%%%%%%%%

Submodular functions arise in:

\[
\boxed{
\text{coverage},
}
\]

\[
\boxed{
\text{information gain},
}
\]

\[
\boxed{
\text{sensor placement},
}
\]

and

\[
\boxed{
\text{influence maximization}.
}
\]

They are often viewed as discrete analogues of convex or concave
structure, depending on whether one studies minimization or maximization.

%%%%%%%%%%%%%%%%%%%%%%%%%%%%%%%%%%%%%%%%%%%%%%%%%%%%%%%%%%%%
\subsection{Monotone Submodular Maximization}
\label{subsec:monotone-submodular-maximization}
%%%%%%%%%%%%%%%%%%%%%%%%%%%%%%%%%%%%%%%%%%%%%%%%%%%%%%%%%%%%

Suppose \(f\) is monotone submodular and one may select at most \(k\)
elements:

\[
\boxed{
\max
f(S)
}
\]

subject to

\[
\boxed{
|S|\leq k.
}
\]

The greedy algorithm repeatedly adds the element with largest marginal
gain.

This greedy method has a strong approximation guarantee.

%%%%%%%%%%%%%%%%%%%%%%%%%%%%%%%%%%%%%%%%%%%%%%%%%%%%%%%%%%%%
\subsection{Submodular Greedy Guarantee}
\label{subsec:submodular-greedy-guarantee}
%%%%%%%%%%%%%%%%%%%%%%%%%%%%%%%%%%%%%%%%%%%%%%%%%%%%%%%%%%%%

For nonnegative monotone submodular maximization under a cardinality
constraint, the standard greedy algorithm achieves at least

\[
\boxed{
1-\frac1e
}
\]

of the optimal value.

Thus,

\[
\boxed{
f(S_{\mathrm{greedy}})
\geq
\left(
1-\frac1e
\right)
f(S^*).
}
\]

This is one of the most important approximation results in discrete
optimization.

%%%%%%%%%%%%%%%%%%%%%%%%%%%%%%%%%%%%%%%%%%%%%%%%%%%%%%%%%%%%
\subsection{Submodular Minimization Preview}
\label{subsec:submodular-minimization-preview}
%%%%%%%%%%%%%%%%%%%%%%%%%%%%%%%%%%%%%%%%%%%%%%%%%%%%%%%%%%%%

Unconstrained submodular minimization can be solved in polynomial time.

This is striking because submodular maximization can be difficult.

Thus the direction of optimization matters greatly even for the same
function class.

%%%%%%%%%%%%%%%%%%%%%%%%%%%%%%%%%%%%%%%%%%%%%%%%%%%%%%%%%%%%
\subsection{Dynamic Programming and State Spaces}
\label{subsec:combinatorial-dp-state-spaces}
%%%%%%%%%%%%%%%%%%%%%%%%%%%%%%%%%%%%%%%%%%%%%%%%%%%%%%%%%%%%

A discrete optimization problem may be represented by states

\[
s
\]

and decisions

\[
a.
\]

The optimal value may satisfy a recurrence such as

\[
\boxed{
V(s)
=
\min_{a\in A(s)}
\left[
c(s,a)
+
V(T(s,a))
\right].
}
\]

This turns combinatorial search into recursive optimization.

The full framework is developed in the next section.

%%%%%%%%%%%%%%%%%%%%%%%%%%%%%%%%%%%%%%%%%%%%%%%%%%%%%%%%%%%%
\subsection{Memoization}
\label{subsec:combinatorial-memoization}
%%%%%%%%%%%%%%%%%%%%%%%%%%%%%%%%%%%%%%%%%%%%%%%%%%%%%%%%%%%%

If many solution paths reach the same subproblem, solving that subproblem
once and storing the answer avoids repeated computation.

This is memoization.

It can reduce exponential recursive enumeration to a much smaller state
space when repeated subproblems are abundant.

%%%%%%%%%%%%%%%%%%%%%%%%%%%%%%%%%%%%%%%%%%%%%%%%%%%%%%%%%%%%
\subsection{Pseudo-Polynomial Dynamic Programming}
\label{subsec:combinatorial-pseudo-polynomial}
%%%%%%%%%%%%%%%%%%%%%%%%%%%%%%%%%%%%%%%%%%%%%%%%%%%%%%%%%%%%

For integer knapsack, a dynamic program can depend on capacity \(W\).

Its complexity may be polynomial in \(W\) but not in the number of bits
required to encode \(W\).

Such an algorithm is pseudo-polynomial rather than polynomial in the
formal input length.

%%%%%%%%%%%%%%%%%%%%%%%%%%%%%%%%%%%%%%%%%%%%%%%%%%%%%%%%%%%%
\subsection{Approximation Schemes Preview}
\label{subsec:approximation-schemes-preview}
%%%%%%%%%%%%%%%%%%%%%%%%%%%%%%%%%%%%%%%%%%%%%%%%%%%%%%%%%%%%

A polynomial-time approximation scheme, abbreviated PTAS, produces for
each fixed

\[
\varepsilon>0
\]

a solution within a factor such as

\[
\boxed{
1+\varepsilon
}
\]

of optimal for a minimization problem, with polynomial running time for
fixed \(\varepsilon\).

%%%%%%%%%%%%%%%%%%%%%%%%%%%%%%%%%%%%%%%%%%%%%%%%%%%%%%%%%%%%
\subsection{FPTAS Preview}
\label{subsec:fptas-preview}
%%%%%%%%%%%%%%%%%%%%%%%%%%%%%%%%%%%%%%%%%%%%%%%%%%%%%%%%%%%%

A fully polynomial-time approximation scheme, abbreviated FPTAS, has
running time polynomial in both:

\[
\boxed{
\text{input size},
}
\]

and

\[
\boxed{
\frac1\varepsilon.
}
\]

The \(0\)-\(1\) knapsack problem admits an FPTAS.

%%%%%%%%%%%%%%%%%%%%%%%%%%%%%%%%%%%%%%%%%%%%%%%%%%%%%%%%%%%%
\subsection{Parameterization Preview}
\label{subsec:parameterized-complexity-preview}
%%%%%%%%%%%%%%%%%%%%%%%%%%%%%%%%%%%%%%%%%%%%%%%%%%%%%%%%%%%%

Some hard problems become manageable when a structural parameter \(k\) is
small.

A fixed-parameter tractable algorithm has running time of the form

\[
\boxed{
f(k)n^{O(1)}.
}
\]

This separates the potentially expensive dependence on \(k\) from the
polynomial dependence on the overall input size.

%%%%%%%%%%%%%%%%%%%%%%%%%%%%%%%%%%%%%%%%%%%%%%%%%%%%%%%%%%%%
\subsection{Treewidth Preview}
\label{subsec:treewidth-preview}
%%%%%%%%%%%%%%%%%%%%%%%%%%%%%%%%%%%%%%%%%%%%%%%%%%%%%%%%%%%%

Treewidth measures how close a graph is to a tree in a structural sense.

Many graph problems that are difficult in general become tractable on
graphs of bounded treewidth using dynamic programming over a tree
decomposition.

This is an important example of exploiting combinatorial structure.

%%%%%%%%%%%%%%%%%%%%%%%%%%%%%%%%%%%%%%%%%%%%%%%%%%%%%%%%%%%%
\subsection{Decomposition}
\label{subsec:combinatorial-decomposition}
%%%%%%%%%%%%%%%%%%%%%%%%%%%%%%%%%%%%%%%%%%%%%%%%%%%%%%%%%%%%

Large problems may contain weakly coupled components.

If the problem separates into parts, one can solve smaller subproblems and
combine their information.

Decomposition ideas include:

\[
\boxed{
\text{dynamic programming},
}
\]

\[
\boxed{
\text{Lagrangian relaxation},
}
\]

\[
\boxed{
\text{Benders decomposition},
}
\]

and

\[
\boxed{
\text{column generation}.
}
\]

%%%%%%%%%%%%%%%%%%%%%%%%%%%%%%%%%%%%%%%%%%%%%%%%%%%%%%%%%%%%
\subsection{Column Generation and Combinatorial Subproblems}
\label{subsec:column-generation-combinatorial}
%%%%%%%%%%%%%%%%%%%%%%%%%%%%%%%%%%%%%%%%%%%%%%%%%%%%%%%%%%%%

In column generation, the master problem may contain an enormous family of
possible combinatorial objects.

For example, one variable may represent:

\[
\boxed{
\text{an entire route},
}
\]

\[
\boxed{
\text{a schedule},
}
\]

or

\[
\boxed{
\text{a cutting pattern}.
}
\]

A pricing problem searches for a new object with attractive reduced cost.

%%%%%%%%%%%%%%%%%%%%%%%%%%%%%%%%%%%%%%%%%%%%%%%%%%%%%%%%%%%%
\subsection{Combinatorial Optimization Workflow}
\label{subsec:combinatorial-optimization-workflow}
%%%%%%%%%%%%%%%%%%%%%%%%%%%%%%%%%%%%%%%%%%%%%%%%%%%%%%%%%%%%

A practical workflow is:

\begin{enumerate}
    \item identify the discrete object to be selected;
    \item determine whether the object is a subset, path, tree, matching,
          permutation, partition, or another structure;
    \item define the objective;
    \item write feasibility rules clearly;
    \item search for known special structure;
    \item determine whether a polynomial-time exact algorithm is known;
    \item formulate a binary or integer model when useful;
    \item examine LP or convex relaxations;
    \item strengthen the relaxation using valid inequalities;
    \item consider dynamic programming for recursive structure;
    \item consider branch-and-bound for exact general search;
    \item consider approximation algorithms for NP-hard problems;
    \item consider heuristics when exact methods are too expensive;
    \item use bounds to evaluate solution quality;
    \item exploit symmetry and decomposition;
    \item verify feasibility;
    \item distinguish empirical quality from proven guarantees.
\end{enumerate}

%%%%%%%%%%%%%%%%%%%%%%%%%%%%%%%%%%%%%%%%%%%%%%%%%%%%%%%%%%%%
\subsection{Choosing a Combinatorial Method}
\label{subsec:choosing-combinatorial-method}
%%%%%%%%%%%%%%%%%%%%%%%%%%%%%%%%%%%%%%%%%%%%%%%%%%%%%%%%%%%%

For shortest paths, consider:

\[
\boxed{
\text{specialized graph algorithms}.
}
\]

For minimum spanning trees, use:

\[
\boxed{
\text{Kruskal or Prim}.
}
\]

For bipartite assignment, use:

\[
\boxed{
\text{matching or assignment algorithms}.
}
\]

For recursive problems with manageable state spaces, use:

\[
\boxed{
\text{dynamic programming}.
}
\]

For general integer formulations, use:

\[
\boxed{
\text{branch-and-cut}.
}
\]

For large NP-hard problems where exact optimization is impractical,
consider:

\[
\boxed{
\text{approximation algorithms or heuristics}.
}
\]

%%%%%%%%%%%%%%%%%%%%%%%%%%%%%%%%%%%%%%%%%%%%%%%%%%%%%%%%%%%%
\subsection{Common Errors}
\label{subsec:combinatorial-optimization-common-errors}
%%%%%%%%%%%%%%%%%%%%%%%%%%%%%%%%%%%%%%%%%%%%%%%%%%%%%%%%%%%%

\subsubsection{Assuming Every Discrete Problem Requires Enumeration}

Many combinatorial problems have specialized polynomial-time algorithms.

Always identify structure before using brute force.

%%%%%%%%%%%%%%%%%%%%%%%%%%%%%%%%%%%%%%%%%%%%%%%%%%%%%%%%%%%%
\subsubsection{Assuming Greedy Algorithms Are Automatically Optimal}

Greedy choice is not a proof.

Optimality requires structure such as matroid exchange properties or an
appropriate graph theorem.

%%%%%%%%%%%%%%%%%%%%%%%%%%%%%%%%%%%%%%%%%%%%%%%%%%%%%%%%%%%%
\subsubsection{Confusing a Heuristic with an Approximation Algorithm}

A heuristic may work well experimentally.

An approximation algorithm additionally has a proven worst-case
performance guarantee.

%%%%%%%%%%%%%%%%%%%%%%%%%%%%%%%%%%%%%%%%%%%%%%%%%%%%%%%%%%%%
\subsubsection{Ignoring the Relaxation Direction}

For minimization, enlarging the feasible set gives a lower bound.

For maximization, enlarging the feasible set gives an upper bound.

The bound direction must be interpreted correctly.

%%%%%%%%%%%%%%%%%%%%%%%%%%%%%%%%%%%%%%%%%%%%%%%%%%%%%%%%%%%%
\subsubsection{Assuming Fractional LP Solutions Are Useless}

Fractional solutions can provide:

\[
\boxed{
\text{bounds},
}
\]

\[
\boxed{
\text{branching information},
}
\]

and

\[
\boxed{
\text{cut-generation information}.
}
\]

They are central to modern discrete optimization.

%%%%%%%%%%%%%%%%%%%%%%%%%%%%%%%%%%%%%%%%%%%%%%%%%%%%%%%%%%%%
\subsubsection{Using Degree Constraints Alone for TSP}

Degree constraints can produce multiple disconnected subtours.

Connectivity or subtour-elimination conditions are also required.

%%%%%%%%%%%%%%%%%%%%%%%%%%%%%%%%%%%%%%%%%%%%%%%%%%%%%%%%%%%%
\subsubsection{Confusing Paths, Trails, and Cycles}

These graph structures have different restrictions.

An optimization formulation must model the intended object precisely.

%%%%%%%%%%%%%%%%%%%%%%%%%%%%%%%%%%%%%%%%%%%%%%%%%%%%%%%%%%%%
\subsubsection{Assuming All Graph Problems Have Similar Complexity}

Shortest path and minimum spanning tree are polynomially solvable.

Maximum clique and general TSP are NP-hard.

Superficially similar graph problems can have very different
computational complexity.

%%%%%%%%%%%%%%%%%%%%%%%%%%%%%%%%%%%%%%%%%%%%%%%%%%%%%%%%%%%%
\subsubsection{Confusing NP-Hard with Impossible}

NP-hard problems can often be solved exactly for moderate instances.

The classification concerns worst-case complexity, not impossibility.

%%%%%%%%%%%%%%%%%%%%%%%%%%%%%%%%%%%%%%%%%%%%%%%%%%%%%%%%%%%%
\subsubsection{Assuming NP-Hard Means Exponential Time Has Been Proven Necessary}

NP-hardness does not currently prove that every exact algorithm must take
exponential time.

Such a conclusion would depend on unresolved complexity assumptions.

%%%%%%%%%%%%%%%%%%%%%%%%%%%%%%%%%%%%%%%%%%%%%%%%%%%%%%%%%%%%
\subsubsection{Using Approximation Ratios Without Stating the Convention}

For minimization and maximization, approximation guarantees are often
written differently.

Always state what the ratio means.

%%%%%%%%%%%%%%%%%%%%%%%%%%%%%%%%%%%%%%%%%%%%%%%%%%%%%%%%%%%%
\subsubsection{Applying Metric-TSP Guarantees to Nonmetric Costs}

Algorithms such as double-tree and Christofides rely on the triangle
inequality.

Their guarantees do not automatically hold for arbitrary TSP instances.

%%%%%%%%%%%%%%%%%%%%%%%%%%%%%%%%%%%%%%%%%%%%%%%%%%%%%%%%%%%%
\subsubsection{Assuming Local Search Finds the Global Optimum}

A locally optimal solution depends on the chosen neighborhood.

Different neighborhoods can have very different local optima.

%%%%%%%%%%%%%%%%%%%%%%%%%%%%%%%%%%%%%%%%%%%%%%%%%%%%%%%%%%%%
\subsubsection{Ignoring Symmetry}

Equivalent combinatorial solutions can create enormous redundant search.

Symmetry-breaking or canonical representations may improve algorithms.

%%%%%%%%%%%%%%%%%%%%%%%%%%%%%%%%%%%%%%%%%%%%%%%%%%%%%%%%%%%%
\subsection{Applications}
\label{subsec:combinatorial-optimization-applications}
%%%%%%%%%%%%%%%%%%%%%%%%%%%%%%%%%%%%%%%%%%%%%%%%%%%%%%%%%%%%

\begin{applicationbox}

\textbf{Transportation and Routing.}
Shortest paths, vehicle routing, TSP variants, and network design are
central combinatorial optimization problems.

\medskip

\textbf{Scheduling.}
Jobs must be assigned, ordered, and timed while satisfying resource and
precedence restrictions.

\medskip

\textbf{Telecommunications.}
Network topology, routing, frequency assignment, and infrastructure
placement involve discrete choices.

\medskip

\textbf{Computer Science.}
Graph algorithms, compiler optimization, memory allocation, clustering,
and data structures often lead to combinatorial optimization.

\medskip

\textbf{Logistics.}
Warehouse selection, fleet assignment, packing, and distribution planning
combine discrete structures with cost optimization.

\medskip

\textbf{Biology.}
Sequence alignment, phylogenetic trees, network reconstruction, and
molecular structure problems contain combinatorial components.

\medskip

\textbf{Machine Learning.}
Feature selection, clustering, subset selection, structured prediction,
and experimental design can involve discrete optimization.

\medskip

\textbf{Public Planning.}
Facility placement, districting, emergency response, and infrastructure
design frequently require optimization over discrete alternatives.

\end{applicationbox}

%%%%%%%%%%%%%%%%%%%%%%%%%%%%%%%%%%%%%%%%%%%%%%%%%%%%%%%%%%%%
\subsection{Exercises}
\label{subsec:combinatorial-optimization-exercises}
%%%%%%%%%%%%%%%%%%%%%%%%%%%%%%%%%%%%%%%%%%%%%%%%%%%%%%%%%%%%

\begin{exercise}[1]
Define combinatorial optimization.
\end{exercise}

\begin{exercise}[1]
Give three examples of discrete feasible structures.
\end{exercise}

\begin{exercise}[1]
How many subsets does an \(n\)-element set have?
\end{exercise}

\begin{exercise}[1]
How many permutations are there of \(n\) distinct objects?
\end{exercise}

\begin{exercise}[1]
Define a matching.
\end{exercise}

\begin{exercise}[1]
Define a spanning tree.
\end{exercise}

\begin{exercise}[1]
Define a vertex cover.
\end{exercise}

\begin{exercise}[1]
Define an independent set.
\end{exercise}

\begin{exercise}[1]
Define a clique.
\end{exercise}

\begin{exercise}[1]
Define an approximation algorithm.
\end{exercise}

\begin{exercise}[1]
Define a submodular function conceptually.
\end{exercise}

\begin{exercise}[2]
Represent a subset of five items using binary variables.
\end{exercise}

\begin{exercise}[2]
Write a constraint selecting exactly three of eight items.
\end{exercise}

\begin{exercise}[2]
Compute the cost of a path whose edge weights are

\[
2,\ 7,\ 4,\ 1.
\]
\end{exercise}

\begin{exercise}[2]
How many edges does a spanning tree on \(12\) vertices contain?
\end{exercise}

\begin{exercise}[2]
Determine whether a given edge set is a matching.
\end{exercise}

\begin{exercise}[2]
Determine whether a given vertex set is independent.
\end{exercise}

\begin{exercise}[2]
Write binary constraints for a vertex-cover problem.
\end{exercise}

\begin{exercise}[2]
Write binary constraints for an independent-set problem.
\end{exercise}

\begin{exercise}[2]
Write a set-cover model for a universe with four elements.
\end{exercise}

\begin{exercise}[2]
Write a set-partitioning constraint requiring every element to be covered
exactly once.
\end{exercise}

\begin{exercise}[3]
Use Dijkstra's algorithm on a small nonnegative weighted graph.
\end{exercise}

\begin{exercise}[3]
Use Kruskal's algorithm to construct a minimum spanning tree.
\end{exercise}

\begin{exercise}[3]
Use Prim's algorithm on the same graph and compare the result.
\end{exercise}

\begin{exercise}[3]
Prove that a tree on \(n\) vertices has \(n-1\) edges.
\end{exercise}

\begin{exercise}[3]
Formulate a weighted bipartite matching problem.
\end{exercise}

\begin{exercise}[3]
Formulate a three-agent, three-task assignment problem.
\end{exercise}

\begin{exercise}[3]
Write degree constraints for a four-city directed TSP.
\end{exercise}

\begin{exercise}[3]
Give an example showing that TSP degree constraints can produce two
subtours.
\end{exercise}

\begin{exercise}[3]
Write a subtour-elimination constraint for a specified subset \(S\).
\end{exercise}

\begin{exercise}[3]
Formulate a minimum vertex-cover problem.
\end{exercise}

\begin{exercise}[3]
Formulate a maximum independent-set problem.
\end{exercise}

\begin{exercise}[3]
Prove that the complement of a vertex cover is an independent set.
\end{exercise}

\begin{exercise}[3]
Prove

\[
\tau(G)+\alpha(G)=|V|.
\]
\end{exercise}

\begin{exercise}[3]
Show that a clique in \(G\) is an independent set in \(\overline{G}\).
\end{exercise}

\begin{exercise}[3]
Formulate graph coloring using binary assignment variables.
\end{exercise}

\begin{exercise}[3]
Construct a maximum-cut formulation for a small graph.
\end{exercise}

\begin{exercise}[3]
Give an example where greedy \(0\)-\(1\) knapsack by value-to-weight ratio
is not optimal.
\end{exercise}

\begin{exercise}[3]
Apply the matroid greedy algorithm to a weighted graphic matroid.
\end{exercise}

\begin{exercise}[3]
Explain why Kruskal's algorithm is a matroid greedy algorithm.
\end{exercise}

\begin{exercise}[3]
Construct an LP relaxation for a combinatorial problem.
\end{exercise}

\begin{exercise}[3]
Explain whether its relaxation gives an upper or lower bound.
\end{exercise}

\begin{exercise}[3]
Compare an integer feasible solution with the relaxation bound.
\end{exercise}

\begin{exercise}[3]
Explain how branch-and-bound can solve a combinatorial problem exactly.
\end{exercise}

\begin{exercise}[3]
Construct a valid inequality that removes a given fractional point.
\end{exercise}

\begin{exercise}[3]
Explain the difference between an integer hull and an LP relaxation.
\end{exercise}

\begin{exercise}[3]
Explain why integrality of a polyhedron is algorithmically valuable.
\end{exercise}

\begin{exercise}[3]
Classify shortest path, MST, and TSP according to known computational
difficulty.
\end{exercise}

\begin{exercise}[3]
Write the decision version of the traveling salesperson problem.
\end{exercise}

\begin{exercise}[3]
Explain the difference between \(P\), \(NP\), NP-complete, and NP-hard.
\end{exercise}

\begin{exercise}[3]
Show why a proposed TSP tour can be verified in polynomial time.
\end{exercise}

\begin{exercise}[3]
Apply the greedy set-cover algorithm to a small instance.
\end{exercise}

\begin{exercise}[3]
Compute

\[
H_5
=
1+\frac12+\frac13+\frac14+\frac15.
\]
\end{exercise}

\begin{exercise}[3]
Explain why a minimum spanning tree provides a lower bound on the metric
TSP optimum.
\end{exercise}

\begin{exercise}[3]
Derive the factor-\(2\) guarantee of the double-tree metric-TSP method.
\end{exercise}

\begin{exercise}[3]
Apply the maximal-matching \(2\)-approximation for vertex cover.
\end{exercise}

\begin{exercise}[3]
Define a \(2\)-opt neighborhood for TSP.
\end{exercise}

\begin{exercise}[3]
Give an example of a local optimum that is not globally optimal.
\end{exercise}

\begin{exercise}[3]
Verify submodularity of a simple coverage function.
\end{exercise}

\begin{exercise}[3]
Apply greedy cardinality-constrained submodular maximization to a small
coverage problem.
\end{exercise}

\begin{exercise}[3]
Explain the difference between polynomial and pseudo-polynomial running
time.
\end{exercise}

\begin{exercise}[3]
Explain what a PTAS guarantees.
\end{exercise}

\begin{exercise}[3]
Explain what an FPTAS guarantees.
\end{exercise}

\begin{exercise}[4]
Study polyhedral descriptions of the spanning-tree polytope.
\end{exercise}

\begin{exercise}[4]
Study the bipartite matching polytope.
\end{exercise}

\begin{exercise}[4]
Prove integrality of the bipartite matching formulation.
\end{exercise}

\begin{exercise}[4]
Study Edmonds' matching algorithm for general graphs.
\end{exercise}

\begin{exercise}[4]
Study the minimum-cost perfect matching problem.
\end{exercise}

\begin{exercise}[4]
Study TSP subtour-elimination polytopes.
\end{exercise}

\begin{exercise}[4]
Investigate comb inequalities for the TSP.
\end{exercise}

\begin{exercise}[4]
Study branch-and-cut methods for the TSP.
\end{exercise}

\begin{exercise}[4]
Study the set-cover polyhedron.
\end{exercise}

\begin{exercise}[4]
Prove the greedy logarithmic approximation guarantee for set cover.
\end{exercise}

\begin{exercise}[4]
Prove the \(2\)-approximation guarantee for vertex cover.
\end{exercise}

\begin{exercise}[4]
Study Christofides' algorithm and prove its
\(\frac32\)-approximation guarantee.
\end{exercise}

\begin{exercise}[4]
Study approximation hardness for general TSP.
\end{exercise}

\begin{exercise}[4]
Study semidefinite relaxations for maximum cut.
\end{exercise}

\begin{exercise}[4]
Study randomized rounding for a covering problem.
\end{exercise}

\begin{exercise}[4]
Study concentration inequalities used in randomized rounding.
\end{exercise}

\begin{exercise}[4]
Prove the matroid greedy theorem.
\end{exercise}

\begin{exercise}[4]
Study matroid intersection.
\end{exercise}

\begin{exercise}[4]
Study weighted matroid intersection.
\end{exercise}

\begin{exercise}[4]
Study submodular minimization.
\end{exercise}

\begin{exercise}[4]
Prove the
\[
1-\frac1e
\]
guarantee for greedy monotone submodular maximization under a cardinality
constraint.
\end{exercise}

\begin{exercise}[4]
Study approximation algorithms under matroid constraints.
\end{exercise}

\begin{exercise}[4]
Study fixed-parameter algorithms for vertex cover.
\end{exercise}

\begin{exercise}[4]
Study dynamic programming on bounded-treewidth graphs.
\end{exercise}

\begin{exercise}[4]
Study the equivalence between optimization and separation for rational
polyhedra.
\end{exercise}

\begin{exercise}[4]
Investigate column generation for set-partitioning models.
\end{exercise}

\begin{exercise}[4]
Study decomposition methods for large combinatorial optimization models.
\end{exercise}

%%%%%%%%%%%%%%%%%%%%%%%%%%%%%%%%%%%%%%%%%%%%%%%%%%%%%%%%%%%%
\subsection{Conceptual Summary}
\label{subsec:combinatorial-optimization-summary}
%%%%%%%%%%%%%%%%%%%%%%%%%%%%%%%%%%%%%%%%%%%%%%%%%%%%%%%%%%%%

Combinatorial optimization solves problems of the form

\[
\boxed{
\min_{x\in\mathcal{F}}
c(x),
}
\]

where \(\mathcal{F}\) is a discrete family of structures.

Typical objects include

\[
\boxed{
\text{subsets}
+
\text{paths}
+
\text{trees}
+
\text{matchings}
+
\text{tours}.
}
\]

Some important combinatorial problems have polynomial-time exact
algorithms:

\[
\boxed{
\text{shortest path},
}
\]

\[
\boxed{
\text{minimum spanning tree},
}
\]

and

\[
\boxed{
\text{bipartite matching}.
}
\]

Others are NP-hard:

\[
\boxed{
\text{TSP},
}
\]

\[
\boxed{
\text{maximum clique},
}
\]

\[
\boxed{
\text{set cover},
}
\]

and

\[
\boxed{
\text{graph coloring}.
}
\]

Relaxations provide bounds:

\[
\boxed{
\mathcal{F}
\subseteq
\mathcal{R}.
}
\]

Exact algorithms may use:

\[
\boxed{
\text{dynamic programming},
}
\]

\[
\boxed{
\text{branch-and-bound},
}
\]

and

\[
\boxed{
\text{branch-and-cut}.
}
\]

When exact optimization is too expensive, one may use approximation
algorithms with provable performance guarantees or heuristics aimed at
high-quality practical solutions.

Combinatorial optimization therefore combines

\[
\boxed{
\text{discrete mathematics}
+
\text{graph theory}
+
\text{integer programming}
+
\text{algorithms}
+
\text{complexity theory}.
}
\]

%%%%%%%%%%%%%%%%%%%%%%%%%%%%%%%%%%%%%%%%%%%%%%%%%%%%%%%%%%%%
\subsection{Section Connections}
\label{subsec:combinatorial-optimization-connections}
%%%%%%%%%%%%%%%%%%%%%%%%%%%%%%%%%%%%%%%%%%%%%%%%%%%%%%%%%%%%

\chapterconnection{
Combinatorial Optimization builds on the discrete structures developed in
Chapter~6 and the integer-programming framework of
Section~\ref{sec:integer-programming}. Graph algorithms, matroids,
polyhedral formulations, relaxations, approximation methods, and
complexity theory explain why some discrete optimization problems admit
efficient exact solutions while others require exponential search,
approximation, or heuristics. The next section, Dynamic Programming,
develops one of the central exact optimization methods for problems with
sequential decisions, optimal substructure, and repeated subproblems.
Dynamic programming will later connect directly with stochastic
optimization, scheduling, control, and decision theory.
}

%%%%%%%%%%%%%%%%%%%%%%%%%%%%%%%%%%%%%%%%%%%%%%%%%%%%%%%%%%%%
% SECTION SOURCE TRACKING
%%%%%%%%%%%%%%%%%%%%%%%%%%%%%%%%%%%%%%%%%%%%%%%%%%%%%%%%%%%%

%------------------------------------------------------------
% 10.6 Combinatorial Optimization
%------------------------------------------------------------
%
% Source basis:
%   - Standard combinatorial optimization theory.
%   - Standard graph algorithms and discrete optimization methods.
%   - Standard computational complexity and approximation theory.
%
% Used for:
%   - discrete feasible structures;
%   - combinatorial explosion;
%   - binary incidence representations;
%   - weighted graph optimization;
%   - shortest paths;
%   - minimum spanning trees;
%   - Kruskal and Prim algorithms;
%   - greedy algorithms;
%   - matching and assignment;
%   - TSP modeling;
%   - degree constraints and subtours;
%   - subtour-elimination constraints;
%   - set cover, packing, and partitioning;
%   - vertex cover;
%   - independent sets;
%   - cliques;
%   - graph coloring;
%   - maximum cut;
%   - knapsack;
%   - matroids;
%   - matroid greedy optimization;
%   - exact algorithms;
%   - dynamic-programming preview;
%   - branch-and-bound;
%   - relaxations;
%   - polyhedral combinatorics;
%   - facets and separation;
%   - integrality;
%   - total unimodularity;
%   - P, NP, NP-completeness, and NP-hardness;
%   - polynomial reductions;
%   - approximation algorithms;
%   - greedy set-cover approximation;
%   - metric-TSP approximation;
%   - Christofides preview;
%   - vertex-cover approximation;
%   - local search;
%   - metaheuristics;
%   - randomized algorithms and rounding;
%   - submodular functions;
%   - submodular maximization;
%   - approximation schemes;
%   - parameterized complexity;
%   - treewidth;
%   - decomposition and column-generation connections;
%   - applications and exercises.
%
% Notes:
%   - Dynamic Programming is developed next in Section 10.7.
%   - Network Optimization is developed in Section 10.10.
%   - Scheduling Theory is developed in Section 10.12.
%   - Graph-theoretic foundations are developed canonically in Chapter 6.
%
%%%%%%%%%%%%%%%%%%%%%%%%%%%%%%%%%%%%%%%%%%%%%%%%%%%%%%%%%%%%